\section{Limitations}
\textbf{Small, clustered test set.} The test split has 8 patients and \NTest{} criteria that are not independent; differences of a few points are within the bootstrap intervals, and rare classes (\emph{not included}, \emph{excluded}) have a handful of examples or none. Conclusions about those classes are therefore weak.

\textbf{Model scale and hardware.} Models of 0.36--1.7B parameters trained for two to four epochs on 602 examples are far from the capability of GPT-4-class systems; we did not tune LoRA rank, learning rate or epoch count beyond the single configuration per model reported, and we did not run multiple seeds, so run-to-run variance is unknown. Epoch selection used validation macro-F1, which moves by several points between epochs and is itself noisy. The GPU numbers come from a shared Kaggle T4 whose clock and neighbours we do not control.

\textbf{Shortness and style of notes.} The corpus contains short synthetic descriptions; real electronic health records are longer, noisier and would need chunking and retrieval before this prompt format applies. The app rejects prompts above 3\,000 tokens rather than truncating silently.

\textbf{Explanations are extractive.} They restate the label and the cited sentences; they are not free-text clinical reasoning and cannot detect an unsupported citation beyond checking that the ID exists in the note.

\textbf{No patient-level ground truth.} The preliminary outcome is a transparent rule over criterion labels; it has not been validated against any patient-trial label and should not be read as one.

\textbf{Benchmark validity.} Timings come from one laptop CPU with background load kept minimal but not zero. Memory is a sampled RSS proxy. The educational tiled attention is pure PyTorch and says nothing about fused GPU kernels; speculative-decoding speed-ups measured here depend on the CPU cost ratio of the two models and on the very short natural outputs.

\textbf{Safety.} The system is an educational screening prototype. It must not be used for enrollment, treatment or any other clinical decision.

\section{Conclusion}
We built, trained and evaluated an end-to-end criterion-level eligibility screener on the TrialGPT annotation corpus with a strictly patient-grouped split, and implemented the Transformer inference mechanisms ourselves rather than through a library switch. Fine-tuning SmolLM2-1.7B with LoRA lifts test macro-F1 from \GBaseMF{} (zero-shot) to \GTestMF{} and accuracy from \GBaseAcc{} to \GTestAcc{}, with evidence F1 \GTestEvF{}; however, the accuracy gain over the \GMajAcc{} majority baseline is not statistically established, the improvement is confined to inclusion criteria, and for exclusion criteria the model almost always answers \emph{not excluded}, treating missing information as a negative finding in 40 of 42 such cases. The explicit KV cache reproduces uncached greedy output exactly and is \GKVSpeedup$\times$ faster at 512+128 tokens at the price of a cache that grows by 393\,kB per token for the 1.7B model; the tiled online-softmax attention matches the reference to $10^{-6}$ and bounds attention memory linearly in sequence length, but is much slower than fused kernels in pure PyTorch; and our greedy speculative decoder is token-exact against the target and against Hugging Face assisted generation, but with a draft that is as deep as the target it is slower than plain decoding on a GPU, which we explain with the cost ratio. Future work should address the exclusion-criterion collapse (class-balanced or cost-sensitive training, an explicit abstention head, or retrieval of relevant sentences), evaluate several seeds, and pair the target with a shallower draft. The system is an educational prototype and not a tool for clinical decisions.
